\chapter{移位寄存器}

\section{先建立数据流动的思维模型}

回顾D触发器：一个D触发器在时钟上升沿采样D输入，输出到Q。

现在把\textbf{多个D触发器首尾连接}：

\begin{center}
\begin{tikzpicture}
  \node[draw, minimum width=1.2cm, minimum height=1.5cm, label=above:DFF0] (dff0) at (0,0) {};
  \node[draw, minimum width=1.2cm, minimum height=1.5cm, label=above:DFF1] (dff1) at (2,0) {};
  \node[draw, minimum width=1.2cm, minimum height=1.5cm, label=above:DFF2] (dff2) at (4,0) {};
  \node[draw, minimum width=1.2cm, minimum height=1.5cm, label=above:DFF3] (dff3) at (6,0) {};

  \draw[->] (dff0.east) -- (dff1.west) node[midway,above] {Q0=D1};
  \draw[->] (dff1.east) -- (dff2.west) node[midway,above] {Q1=D2};
  \draw[->] (dff2.east) -- (dff3.west) node[midway,above] {Q2=D3};

  \node (din) at (-2,0) {$D_{in}$};
  \draw[->] (din) -- (dff0.west);
  \node (dout) at (8,0) {$Q_3$};
  \draw[->] (dff3.east) -- (dout);

  \node at (3,-1.5) {每个D触发器的Q接到下一个D触发器的D};
\end{tikzpicture}
\end{center}

\begin{keypoint}[移位寄存器的核心结构]
移位寄存器 = N个D触发器串联

$$\text{DFF}_k \text{的} Q \to \text{DFF}_{k+1} \text{的} D$$

在时钟上升沿，每个DFF采样前一级DFF的Q输出。
结果：\textbf{所有数据同时向右移动一位。}
\end{keypoint}

\section{逐拍分析：数据如何移动}

这是全书最重要的状态分析之一。请仔细跟着每一拍。

\subsection{初始条件}

假设4个D触发器初始值：Q3Q2Q1Q0 = 0000

串行输入 $D_{in}$ 序列：1, 0, 1, 1, 0, ...

（每个时钟周期，一个新的$D_{in}$值被送入最左边的DFF0）

\subsection{时钟第0拍（初始状态）}

\begin{itemize}
  \item 寄存器内容：Q3=0, Q2=0, Q1=0, Q0=0 → [0 0 0 0]
  \item 当前$D_{in}$ = 1（准备好，等待下一个上升沿）
\end{itemize}

\subsection{时钟第1拍上升沿}

\begin{beatanalysis}[移位寄存器——时钟第1拍]
\textbf{上升沿到达时发生的事情（所有DFF同时动作）：}

\begin{center}
\begin{tabular}{|c|c|c|c|c|}
\hline
\textbf{DFF} & \textbf{D输入(上升沿前)} & \textbf{上升沿后Q} & \textbf{DFF动作} \\
\hline
DFF0 & $D_{in}$=1 & Q0=\textbf{1} & 采样外部输入1 \\
DFF1 & Q0=0 & Q1=\textbf{0} & 采样前一拍Q0的值(0) \\
DFF2 & Q1=0 & Q2=\textbf{0} & 采样前一拍Q1的值(0) \\
DFF3 & Q2=0 & Q3=\textbf{0} & 采样前一拍Q2的值(0) \\
\hline
\end{tabular}

\textbf{上升沿后寄存器内容：[0 0 0 1]}（1进入最右端）

原来在Q0的0移到了Q1，原来在Q1的0移到了Q2，原来在Q2的0移到了Q3。
Q3原来的值（0）被"挤出去"了——这就是串行输出。
\end{center}
\end{beatanalysis}

\subsection{时钟第2拍上升沿}

现在 $D_{in}$ = 0

\begin{beatanalysis}[移位寄存器——时钟第2拍]
\begin{center}
\begin{tabular}{|c|c|c|c|c|}
\hline
\textbf{DFF} & \textbf{D输入(上升沿前)} & \textbf{上升沿后Q} & \textbf{DFF动作} \\
\hline
DFF0 & $D_{in}$=0 & Q0=\textbf{0} & 采样外部输入0 \\
DFF1 & Q0=1 & Q1=\textbf{1} & 采样Q0的值(1) → 1右移 \\
DFF2 & Q1=0 & Q2=\textbf{0} & 采样Q1的值(0) → 0右移 \\
DFF3 & Q2=0 & Q3=\textbf{0} & 采样Q2的值(0) → 0右移 \\
\hline
\end{tabular}

\textbf{上升沿后寄存器内容：[0 0 1 0]}（之前Q0的1现在移到了Q1）
\end{center}
\end{beatanalysis}

\subsection{时钟第3拍上升沿}

现在 $D_{in}$ = 1

\begin{beatanalysis}[移位寄存器——时钟第3拍]
\begin{center}
\begin{tabular}{|c|c|c|c|c|}
\hline
\textbf{DFF} & \textbf{D输入(上升沿前)} & \textbf{上升沿后Q} & \textbf{说明} \\
\hline
DFF0 & $D_{in}$=1 & Q0=\textbf{1} & 新数据进入 \\
DFF1 & Q0=0 & Q1=\textbf{0} & \\
DFF2 & Q1=1 & Q2=\textbf{1} & 1继续右移 \\
DFF3 & Q2=0 & Q3=\textbf{0} & \\
\hline
\end{tabular}

\textbf{上升沿后寄存器内容：[0 1 0 1]}（可以看到1在向右移动！）
\end{center}
\end{beatanalysis}

\subsection{时钟第4拍上升沿}

现在 $D_{in}$ = 1

\textbf{上升沿后寄存器内容：[1 0 1 1]}

\subsection{完整数据移动轨迹}

\begin{center}
\begin{tabular}{|c|c|c|c|c|c|}
\hline
\textbf{时钟拍} & \textbf{$D_{in}$} & \textbf{Q3} & \textbf{Q2} & \textbf{Q1} & \textbf{Q0} \\
\hline
0（初始）& — & 0 & 0 & 0 & 0 \\
1 & 1 & 0 & 0 & 0 & \textbf{1} \\
2 & 0 & 0 & 0 & \textbf{1} & 0 \\
3 & 1 & 0 & \textbf{1} & 0 & 1 \\
4 & 1 & \textbf{1} & 0 & 1 & 1 \\
5 & 0 & \textbf{0} & 1 & 1 & 0 \\
\hline
\end{tabular}
\end{center}

\textbf{观察数据流：}每个1从$D_{in}$进入，经过4个时钟周期后从Q3串行输出。

就像水管中的水流——从一端进入，经过管道，从另一端流出。

\section{核心概念总结}

\begin{enumerate}
  \item \textbf{串行输入（Serial In）}：数据从最左边的DFF一个bit一个bit地进入
  \item \textbf{串行输出（Serial Out）}：数据从最右边的DFF一个bit一个bit地流出
  \item \textbf{并行输出（Parallel Out）}：所有Q值同时观察——看到"管道"中当前的内容
  \item \textbf{并行输入（Parallel Load）}：直接设置所有DFF的值（跳过移位过程）
\end{enumerate}

\section{单向移位寄存器：VHDL实现}

\begin{lstlisting}[caption={4位右移移位寄存器（串入串出/并出）}]
library ieee;
use ieee.std_logic_1164.all;

entity shift_reg_right is
  port (
    clk     : in  std_logic;
    rst_n   : in  std_logic;
    si      : in  std_logic;  -- 串行输入
    so      : out std_logic;  -- 串行输出
    po      : out std_logic_vector(3 downto 0)  -- 并行输出
  );
end shift_reg_right;

architecture behavioral of shift_reg_right is
  signal reg : std_logic_vector(3 downto 0);
begin
  process(clk, rst_n)
  begin
    if rst_n = '0' then
      reg <= (others => '0');
    elsif rising_edge(clk) then
      reg <= si & reg(3 downto 1);  -- 右移：新数据进最高位
    end shift;
  end process;

  so <= reg(0);  -- 串行输出 = 最低位
  po <= reg;     -- 并行输出 = 整个寄存器
end behavioral;
\end{lstlisting}

\textbf{关键语句}：\texttt{reg <= si \& reg(3 downto 1)}

这表示：
\begin{itemize}
  \item reg(3)（原来的最高位）$\to$ reg(2)
  \item reg(2) $\to$ reg(1)
  \item reg(1) $\to$ reg(0)
  \item si（新数据）$\to$ reg(3)
\end{itemize}

数据向右移动一位。新数据从左边进入，最右边的数据（reg(0)）被挤出。

\section{双向移位寄存器}

双向 = 既可以左移也可以右移，由方向控制信号决定。

\begin{lstlisting}[caption={双向移位寄存器}]
library ieee;
use ieee.std_logic_1164.all;

entity shift_reg_bidir is
  port (
    clk     : in  std_logic;
    rst_n   : in  std_logic;
    dir     : in  std_logic;  -- 方向：'1'=右移, '0'=左移
    si      : in  std_logic;  -- 串行输入
    po      : out std_logic_vector(3 downto 0)
  );
end shift_reg_bidir;

architecture behavioral of shift_reg_bidir is
  signal reg : std_logic_vector(3 downto 0);
begin
  process(clk, rst_n)
  begin
    if rst_n = '0' then
      reg <= (others => '0');
    elsif rising_edge(clk) then
      if dir = '1' then
        reg <= si & reg(3 downto 1);  -- 右移
      else
        reg <= reg(2 downto 0) & si;  -- 左移
      end if;
    end if;
  end process;

  po <= reg;
end behavioral;
\end{lstlisting}

\section{带并行装载的移位寄存器}

增加一个load信号，允许直接设置所有D触发器的值：

\begin{lstlisting}[caption={带并行装载的移位寄存器}]
library ieee;
use ieee.std_logic_1164.all;

entity shift_reg_load is
  port (
    clk     : in  std_logic;
    rst_n   : in  std_logic;
    load    : in  std_logic;  -- 并行装载使能
    data_in : in  std_logic_vector(3 downto 0);  -- 并行输入
    si      : in  std_logic;  -- 串行输入
    po      : out std_logic_vector(3 downto 0)
  );
end shift_reg_load;

architecture behavioral of shift_reg_load is
  signal reg : std_logic_vector(3 downto 0);
begin
  process(clk, rst_n)
  begin
    if rst_n = '0' then
      reg <= (others => '0');
    elsif rising_edge(clk) then
      if load = '1' then
        reg <= data_in;  -- 并行装载
      else
        reg <= si & reg(3 downto 1);  -- 右移
      end if;
    end if;
  end process;

  po <= reg;
end behavioral;
\end{lstlisting}

\section{四种输入输出组合}

\begin{center}
\begin{tabular}{|l|l|l|l|}
\hline
\textbf{类型} & \textbf{输入方式} & \textbf{输出方式} & \textbf{典型应用} \\
\hline
SISO & 串行入 & 串行出 & 延迟线 \\
SIPO & 串行入 & 并行出 & 串并转换 \\
PISO & 并行入 & 串行出 & 并串转换 \\
PIPO & 并行入 & 并行出 & 通用寄存器 \\
\hline
\end{tabular}
\end{center}

SISO = Serial In Serial Out（串入串出）
SIPO = Serial In Parallel Out（串入并出）
PISO = Parallel In Serial Out（并入串出）
PIPO = Parallel In Parallel Out（并入并出）

\section{移位寄存器实现分频器}

\subsection{基本原理}

移位寄存器天然具有分频能力。回顾环形计数器和Johnson计数器：

\begin{itemize}
  \item N位环形计数器：N个状态循环，任何Q输出的频率 = $f_{clk} / N$
  \item N位Johnson计数器：2N个状态循环，任何Q输出的频率 = $f_{clk} / (2N)$
\end{itemize}

因为移位寄存器中的"1"（或状态模式）每N拍（或2N拍）才经过同一个触发器一次，所以每个Q输出的翻转频率恰好是时钟的$1/N$（或$1/2N$）。

\subsection{逐拍分析：4位环形计数器做4分频}

假设4位环形计数器初始状态1000，反馈Q0$\to$D3（闭环）。

\begin{center}
\begin{tabular}{|c|c|c|c|c|c|}
\hline
时钟拍 & Q3 & Q2 & Q1 & Q0 & 说明 \\
\hline
0 & 1 & 0 & 0 & 0 & 初始 \\
1 & 0 & 1 & 0 & 0 & 1右移 \\
2 & 0 & 0 & 1 & 0 & \\
3 & 0 & 0 & 0 & 1 & \\
4 & 1 & 0 & 0 & 0 & 回到初始！\\
\hline
\end{tabular}
\end{center}

Q3的输出序列：1, 0, 0, 0, 1, 0, 0, 0, ...
周期 = 4个时钟 $\to$ 频率 = $f_{clk} / 4$。

\textbf{每个Q输出都是4分频信号，只是相位不同。}

\subsection{Johnson计数器做分频器}

4位Johnson计数器（8状态）：
$$0000\to1000\to1100\to1110\to1111\to0111\to0011\to0001\to0000$$

任意Q输出的周期 = 8个时钟 $\to$ 频率 = $f_{clk} / 8$。

\textbf{规律：}N位Johnson计数器 = $2N$分频（所有Q输出频率相同，相位不同）。

\subsection{VHDL实现：移位寄存器型分频器}

\begin{lstlisting}[caption={4位环形计数器做4分频器}]
library ieee;
use ieee.std_logic_1164.all;

entity shift_div4 is
  port(
    clk, rst : in  std_logic;
    clk_out  : out std_logic
  );
end shift_div4;

architecture behav of shift_div4 is
  signal reg : std_logic_vector(3 downto 0);
begin
  process(clk, rst)
  begin
    if rst = '1' then reg <= "1000";
    elsif rising_edge(clk) then
      reg <= reg(0) & reg(3 downto 1);
    end if;
  end process;
  clk_out <= reg(3);  -- Q3 = 4分频输出
end behav;
\end{lstlisting}

\subsection{与计数器分频的对比}

\begin{center}
\begin{tabular}{|l|c|c|}
\hline
\textbf{分频方式} & \textbf{结构} & \textbf{可得分频比} \\
\hline
二进制计数器取Qn & 加法器+寄存器 & $2^n$（Q0=2, Q1=4, Q2=8...）\\
环形计数器 & N个DFF+反馈 & N（4位$\to$4分频）\\
Johnson计数器 & N个DFF+反相反馈 & 2N（4位$\to$8分频）\\
模M计数器+比较 & 加法器+寄存器+比较器 & 任意M \\
\hline
\end{tabular}
\end{center}

\textbf{移位寄存器分频的优缺点：}
\begin{itemize}
  \item 优点：结构规整（只有D触发器和反馈线），关键路径极短（1级DFF延迟），适合超高速分频
  \item 缺点：分频比固定（N或2N），灵活性不如计数器方案
\end{itemize}

\section{移位寄存器实现序列发生器}

\subsection{方法一：循环移位型（只能循环输出初值，不是真正的序列发生器）}

\textbf{先说清楚：这不是真正的序列发生器。}它只是把装进去的值循环输出——装1001就循环出1001，装0101就循环出0101。你无法用它"设计"出一个新的序列。

那为什么还要讲？因为\textbf{硬件结构}和真正的序列发生器一模一样：都是移位寄存器+反馈线。理解了循环移位，查表法就水到渠成。

\textbf{方案：}移位寄存器 + 反馈逻辑（Q0直接反馈给D3）。每拍数据右移一位，最低位移回最高位。

\textbf{例题：}4位寄存器，反馈 = Q0→D3，初值"1001"。

\begin{center}
\begin{tabular}{|c|c|c|c|c|}
\hline
时钟拍 & Q3 & Q2 & Q1 & Q0 \\
\hline
0 & 1 & 0 & 0 & 1（初始）\\
1 & 1 & 1 & 0 & 0（Q0=1反馈到D3）\\
2 & 0 & 1 & 1 & 0 \\
3 & 0 & 0 & 1 & 1 \\
4 & 1 & 0 & 0 & 1（回到初始！）\\
\hline
\end{tabular}
\end{center}

序列长度 = 4拍。输出就是初值本身，循环重复。\textbf{这不是序列"生成"，只是"循环播放"。}

\textbf{【VHDL——循环右移序列发生器】}
\begin{lstlisting}
library ieee;
use ieee.std_logic_1164.all;

entity shift_seq_gen is
  port(clk, rst : in std_logic; seq_out : out std_logic_vector(3 downto 0));
end shift_seq_gen;

architecture behav of shift_seq_gen is
  signal reg : std_logic_vector(3 downto 0);
begin
  process(clk, rst)
  begin
    if rst='1' then reg <= "1001";
    elsif rising_edge(clk) then
      reg <= reg(0) & reg(3 downto 1);  -- 循环右移
    end if;
  end process;
  seq_out <= reg;
end behav;
\end{lstlisting}

\subsection{方法二：查表反馈型（通用序列发生器）}

循环移位只能产生移位相关的序列。要产生\textbf{任意序列}，需要用查表法（CASE语句）根据当前状态决定反馈值。

\textbf{核心结构：}移位寄存器 + CASE查表器。CASE根据寄存器的当前值，查表输出下一位要移入的din值。

\subsection{例题：用3位移位寄存器产生序列 00011101 00011101...}

\textbf{设计分析：}3位移位寄存器共8个状态，每个状态通过CASE指定反馈位din的值。输出取寄存器最高位q(2)。

\subsection{核心问题：CASE表是怎么得到的？}

这是初学者最容易困惑的地方。答案：\textbf{用3位滑动窗口在目标序列上滑动}。

目标序列是：0 0 0 1 1 1 0 1（周期8，循环重复）。

\begin{center}
\begin{tikzpicture}[scale=0.9]
% Target sequence
\node at (0, 2.5) {目标序列(循环):};
\foreach \i/\b in {0/0,1/0,2/0,3/1,4/1,5/1,6/0,7/1} {
  \node at (\i*1.2, 2) {\b};
}
% Window 0
\draw[red,very thick] (-0.3,1.5) rectangle (2.1,1.0);
\node[red] at (0.9, 1.25) {窗口0: 000};
\node[red] at (3.3, 1.25) {$\to$ 下一个是1 $\to$ din='1'};
\node[red] at (3.3, 0.75) {$\to$ state="000"};

% Window 1
\draw[blue,very thick] (0.9,1.5) rectangle (3.3,1.0);
\node[blue] at (2.1, 0.5) {窗口1: 001 $\to$ 下一个1 $\to$ din='1'};

% Window 2
\draw[green!50!black,very thick] (2.1,1.5) rectangle (4.5,1.0);
\node[green!50!black] at (3.3, 0.5) {窗口2: 011 $\to$ 下一个1 $\to$ din='1'};

% Window 3
\draw[orange,very thick] (3.3,1.5) rectangle (5.7,1.0);
\node[orange] at (4.5, 0.5) {窗口3: 111 $\to$ 下一个0 $\to$ din='0'};

% Ellipsis
\node at (6.5, 0.5) {...};
\end{tikzpicture}
\end{center}

\textbf{推导步骤：}

\begin{enumerate}
\item 把目标序列写成循环：\textbf{0 0 0 1 1 1 0 1 | 0 0 0 1 1 1 0 1 | ...}（首尾相接，无限重复）
\item 用宽度为3的窗口，从序列开头每次移动一位，读取窗口内的3个bit——这就是3个D触发器的\textbf{当前状态}
\item 窗口右边紧挨着的下一个bit——就是需要反馈进入的\textbf{din值}
\item 覆盖所有8种可能的3-bit组合（因为3位有$2^3=8$种组合，恰好一个周期遍历完）
\end{enumerate}

\textbf{图解推导——3位滑动窗口法：}

\begin{center}
\begin{tabular}{cccccccccccccccccccc}
\multicolumn{20}{c}{目标序列（无限循环，用竖线标出8个3位窗口）} \\
& & & \fbox{0 0 0} & 1 & 1 & 1 & 0 & 1 & 0 & 0 & 0 & 1 & ... & $\leftarrow$ 窗口0: state=000, din=1 \\
& & 0 & \fbox{0 0 1} & 1 & 1 & 0 & 1 & 0 & 0 & 0 & 1 & ... & & $\leftarrow$ 窗口1: state=001, din=1 \\
& 0 & 0 & \fbox{0 1 1} & 1 & 0 & 1 & 0 & 0 & 0 & 1 & ... & & & $\leftarrow$ 窗口2: state=011, din=1 \\
0 & 0 & 0 & \fbox{1 1 1} & 0 & 1 & 0 & 0 & 0 & 1 & ... & & & & $\leftarrow$ 窗口3: state=111, din=0 \\
0 & 0 & 0 & 1 & \fbox{1 1 0} & 1 & 0 & 0 & 0 & 1 & ... & & & & $\leftarrow$ 窗口4: state=110, din=1 \\
0 & 0 & 0 & 1 & 1 & \fbox{1 0 1} & 0 & 0 & 0 & 1 & ... & & & & $\leftarrow$ 窗口5: state=101, din=0 \\
0 & 0 & 0 & 1 & 1 & 1 & \fbox{0 1 0} & 0 & 0 & 1 & ... & & & & $\leftarrow$ 窗口6: state=010, din=0 \\
0 & 0 & 0 & 1 & 1 & 1 & 0 & \fbox{1 0 0} & 0 & 1 & ... & & & & $\leftarrow$ 窗口7: state=100, din=0 \\
\end{tabular}
\end{center}

\textbf{关键理解：}
\begin{itemize}
  \item \fbox{框内3位} = 3个D触发器的当前状态 (qin)
  \item 框右边紧跟着的bit = 下一个要移入的din值
  \item 框最左边的bit = 当前输出值 (Q2)
  \item 下一个时钟沿：din从右边移入，所有bit左移一位 → 窗口向右滑动一位 → Q2输出下一个bit
\end{itemize}

\subsection{关键技巧：需要几个D触发器？——试3个不行就换4个}

\textbf{窗口宽度 = 寄存器的个数。}窗口太窄，不同位置的窗口值会重复（"撞车"），导致CASE表出现矛盾——同一个状态对应两个不同的din值。

\textbf{举例：序列 00110011...（8位循环），先试3位窗口}

\[
\underbrace{001}_{\text{窗口0}}10011\underbrace{001}_{\text{窗口4}}10011...
\]

\begin{center}
\begin{tabular}{c|c|c}
窗口位置 & 窗口值(3位) & 下一个bit(din) \\
\hline
0 & 001 & 1 \\
1 & 011 & 0 \\
2 & 110 & 0 \\
3 & 100 & 1 \\
4 & 001 & \textbf{？} ← 窗口值=001又出现了！但这次下一个bit是1？0？\\
\end{tabular}
\end{center}

窗口0（001）的下一个是1，窗口4（001）的下一个本来是接着序列...但\textbf{同一个状态"001"在CASE表中只能有一行}，不能同时对应din=1和din=0。这就是"撞车"。

\textbf{解决方案：窗口加宽→换4个寄存器}

\begin{center}
\begin{tabular}{c|c|c}
窗口位置 & 窗口值(4位) & 下一个bit(din) \\
\hline
0 & 0011 & 0 \\
1 & 0110 & 0 \\
2 & 1100 & 1 \\
3 & 1001 & 1 \\
4 & 0011 & 0 ← 和窗口0一样！但下一个bit也一样(都是0)，不冲突 ✓ \\
\end{tabular}
\end{center}

4位窗口下，0011→下一个是0，再次出现0011时下一个还是0，\textbf{CASE表不矛盾}。

\textbf{通用法则：}

\begin{keypoint}[寄存器位数的确定]
\textbf{先试$k = \lceil\log_2 L\rceil$位（L=序列长度），用滑动窗口检查有没有"撞车"。}

\begin{itemize}
  \item 8种窗口值全不同 → $k$位够用
  \item 有重复但重复处的din也相同 → $k$位够用
  \item 有重复且din不同 → \textbf{窗口加宽为$k+1$位，重新检查}
\end{itemize}

选最小不撞车的位数作为寄存器个数。
\end{keypoint}

\subsection{完整例题：你需要几个寄存器？}

\textbf{题目：用查表法产生序列 0011001 0011001...（7位循环），最少需要几个D触发器？}

\textbf{第一步：算起步位数} $\lceil\log_2 7\rceil = 3$位。\textbf{先用3位窗口试。}

把序列循环展开写两遍，方便看环绕：0 0 1 1 0 0 1 | 0 0 1 1 0 0 1

\begin{center}
\begin{tabular}{c|c|c}
位置 & 3位窗口值(qin) & 右边一个bit(din) \\
\hline
0 & 001 & \textbf{1} \\
1 & 011 & \textbf{0} \\
2 & 110 & \textbf{0} \\
3 & 100 & \textbf{0} \\
4 & 001 & \textbf{0} ？\\
5 & 010 & \textbf{0} \\
6 & 100 & \textbf{1} ？\\
\hline
\end{tabular}
\end{center}

\textbf{查撞车：}
\begin{itemize}
  \item 位置0的001 → din=1，位置4的001 → din=0。\textbf{同一个001，din不同！撞车！}
  \item 位置3的100 → din=0，位置6的100 → din=1。\textbf{又撞车！}
\end{itemize}

3位不行→\textbf{换4位窗口。}

\textbf{第二步：4位窗口}

序列：0 0 1 1 0 0 1 | 0 0 1 1 0 0 1

\begin{center}
\begin{tabular}{c|c|c}
位置 & 4位窗口值(qin) & 下一个bit(din) \\
\hline
0 & 0011 & 0 \\
1 & 0110 & 0 \\
2 & 1100 & 1 \\
3 & 1001 & 0 \\
4 & 0010 & 0 \\
5 & 0100 & 1 \\
6 & 1001 & \textbf{1} ？\\
\hline
\end{tabular}
\end{center}

位置3(1001)→din=0，位置6(1001)→din=1。\textbf{还是撞！}→ 换5位。

\textbf{第三步：5位窗口}

\begin{center}
\begin{tabular}{c|c|c}
位置 & 5位窗口值(qin) & 下一个bit(din) \\
\hline
0 & 00110 & 0 \\
1 & 01100 & 1 \\
2 & 11001 & 0 \\
3 & 10010 & 0 \\
4 & 00100 & 1 \\
5 & 01001 & 0 \\
6 & 10011 & (回0)0 \\
\hline
\end{tabular}
\end{center}

\textbf{7个5位窗口全部不同！}→ 需要\textbf{5个寄存器}。

\textbf{结论：}$\lceil\log_2 7\rceil = 3$只是理论下界。实际寄存器数取决于序列本身的结构——有些序列3位够用（如00011101），有些需要更多（如0011001需要5位）。\textbf{试到不撞车为止。}

\subsection{速查：多少位寄存器够用？}

\begin{center}
\begin{tabular}{|c|c|c|c|}
\hline
序列长度L & $\lceil\log_2 L\rceil$（起步） & 最坏需要 & 常见需要 \\
\hline
4 & 2 & 2 & 2 \\
5 & 3 & 3 & 3 \\
6 & 3 & 3$\sim$4 & 3 \\
7 & 3 & 5 & 3$\sim$4 \\
8 & 3 & 3 & 3 \\
\hline
\end{tabular}
\end{center}

\textbf{考试做法：}先用$\lceil\log_2 L\rceil$位，画滑动窗口表，检查有无撞车（同一窗口值对应不同din）。不撞→够用；撞→加一位重新画。

\subsection{VHDL实现：5位寄存器产生序列0011001}

5位寄存器共$2^5=32$种可能状态，但循环只经过7个。CASE只写这7个合法状态，其余25个状态统一扔到when others回初始态。

\begin{lstlisting}
library IEEE;
use IEEE.STD_LOGIC_1164.all;

entity seqgen_5bit is
  port(clk, rst : in std_logic; q : out std_logic);
end seqgen_5bit;

architecture behav of seqgen_5bit is
  signal qin : std_logic_vector(4 downto 0);
  signal din : std_logic;
begin
  -- 查表产生反馈位din
  process(clk)
  begin
    if rising_edge(clk) then
      case qin is
        when "00110" => din <= '0';
        when "01100" => din <= '1';
        when "11001" => din <= '0';
        when "10010" => din <= '0';
        when "00100" => din <= '1';
        when "01001" => din <= '0';
        when "10011" => din <= '0';
        when others  => din <= '0';  -- 非法状态兜底
      end case;
    end if;
  end process;

  -- 移位寄存器：左移，din进入最低位
  process(clk, rst)
  begin
    if rst = '1' then
      qin <= "00110";  -- 初始窗口0
    elsif rising_edge(clk) then
      qin <= qin(3 downto 0) & din;
    end if;
  end process;

  q <= qin(4);  -- 输出最高位
end behav;
\end{lstlisting}

\textbf{逐拍验证：}
\begin{center}
\begin{tabular}{|c|c|c|c|}
\hline
拍 & qin(当前) & din & q(输出)=qin(4) \\
\hline
0 & 00110 & 0 & 0 \\
1 & 01100 & 1 & 0 \\
2 & 11001 & 0 & 1 \\
3 & 10010 & 0 & 1 \\
4 & 00100 & 1 & 0 \\
5 & 01001 & 0 & 0 \\
6 & 10011 & 0 & 1 \\
7 & 00110 & — & 0 ← 回到拍0 \\
\hline
\end{tabular}
\end{center}

输出 q：0,0,1,1,0,0,1,0,0,1,1,0,0,1,... ← 确实就是0011001循环。

\textbf{这就是CASE表的来源：}不撞车的窗口值→下一个bit，整理成表。

\textbf{这就是CASE表的来源：}把上面8个窗口的"窗口值→下一个bit"整理成表，就是CASE表。

\textbf{完整CASE表推导：}
\begin{center}
\begin{tabular}{|c|c|c|c|c|}
\hline
窗口编号 & 窗口内容(qin) & 窗口右边bit(din) & 下一拍qin(左移后) & 输出q(2) \\
\hline
0 & 000 & \textbf{1} & 001 & 0 \\
1 & 001 & \textbf{1} & 011 & 0 \\
2 & 011 & \textbf{1} & 111 & 0 \\
3 & 111 & \textbf{0} & 110 & 1 \\
4 & 110 & \textbf{1} & 101 & 1 \\
5 & 101 & \textbf{0} & 010 & 1 \\
6 & 010 & \textbf{0} & 100 & 0 \\
7 & 100 & \textbf{0} & 000 & 1 \\
\hline
\end{tabular}
\end{center}

\textbf{逐拍追踪：}
\begin{center}
\begin{tabular}{|c|c|c|c|c|}
\hline
时钟拍 & qin(当前) & CASE查表→din & 下一拍qin & 输出q(2) \\
\hline
0 & 000 & 1 & 001 & 0 \\
1 & 001 & 1 & 011 & 0 \\
2 & 011 & 1 & 111 & 0 \\
3 & 111 & 0 & 110 & 1 \\
4 & 110 & 1 & 101 & 1 \\
5 & 101 & 0 & 010 & 1 \\
6 & 010 & 0 & 100 & 0 \\
7 & 100 & 0 & 000 & 1 \\
8 & 000 & — & — & 回到拍0 \\
\hline
\end{tabular}
\end{center}

输出序列 q(2)：0,0,0,1,1,1,0,1,0,0,0,1,1,1,0,1,... 周期=8。\textbf{与目标序列完全一致！}

\textbf{【VHDL——查表反馈型序列发生器】}
\begin{lstlisting}
library IEEE;
use IEEE.STD_LOGIC_1164.all;

entity seqgen_table is
  port(clk, rst : in std_logic; q : out std_logic);
end seqgen_table;

architecture behav of seqgen_table is
  signal qin : std_logic_vector(2 downto 0);
  signal din : std_logic;
begin
  -- 第一段：CASE查表产生反馈位din
  comb_proc: process(clk)
  begin
    if rising_edge(clk) then
      case qin is
        when "000" => din <= '1';
        when "001" => din <= '1';
        when "011" => din <= '1';
        when "111" => din <= '0';
        when "110" => din <= '1';
        when "101" => din <= '0';
        when "010" => din <= '0';
        when "100" => din <= '0';
        when others => din <= '0';
      end case;
    end if;
  end process;

  -- 第二段：移位寄存器（左移，din进入最低位）
  shift_proc: process(clk, rst)
  begin
    if rst = '1' then qin <= (others => '0');
    elsif rising_edge(clk) then
      qin <= qin(1 downto 0) & din;
    end if;
  end process;

  q <= qin(2);  -- 输出最高位
end behav;
\end{lstlisting}

\textbf{【设计思想】}查表法的本质：把想要的序列"写死"在CASE表里。当前状态→查表→下一个输入位。只要改CASE表的内容，就能生成任何8位循环序列。这是\textbf{通用序列发生器}的设计范式。

\subsection{两种方案对比}

\begin{center}
\begin{tabular}{|l|c|c|}
\hline
\textbf{方案} & \textbf{循环移位型} & \textbf{查表反馈型} \\
\hline
反馈逻辑 & reg(0) $\to$ D(N-1)（一根线） & CASE表（多条件判断） \\
序列来源 & 就是初值本身（装入什么输出什么）& CASE表任意指定（真正的设计）\\
代码复杂度 & 1行 & N行（每个状态1行）\\
修改序列 & 改初值（但只能输出初值循环）& 修改CASE表内容（任意序列）\\
适用场景 & 学习移位寄存器工作原理 & 实际的序列发生器设计 \\
\hline
\end{tabular}
\end{center}

\section{移位寄存器实现串行加法器}

\subsection{原理}

并行加法器（如行波进位）一次完成所有位的加法。串行加法器逐位相加，每次处理1位，N拍后得到N位结果。\textbf{核心结构：}1位全加器 + 进位D触发器 + 两个移位寄存器（存放操作数A和B）+ 一个移位寄存器（存放结果）。

\subsection{逐拍工作过程}

以4位加法 1011 + 0110 为例：

\begin{center}
\begin{tabular}{|c|c|c|c|c|c|}
\hline
拍 & A移位器(Q0) & B移位器(Q0) & Cin(DFF) & Sum位 & Cout \\
\hline
0（初）& 1 & 0 & 0 & — & — \\
1 & 1 & 0 & 0 & 1 & 0 \\
2 & 1 & 1 & 0 & 0 & 1 \\
3 & 0 & 1 & 1 & 0 & 1 \\
4 & 1 & 0 & 1 & 0 & 1 \\
\hline
\end{tabular}
\end{center}

每拍：A和B各右移1位（最低位进入全加器），全加器计算Sum和Cout，Cout存入进位DFF作为下一拍的Cin。结果Sum逐位存入结果移位寄存器。4拍后结果移位寄存器 = 0001 0001（即17，1011+0110=17）。

\textbf{【VHDL——4位串行加法器】}
\begin{lstlisting}
library ieee;
use ieee.std_logic_1164.all;

entity serial_adder is
  port(
    clk, rst, start : in std_logic;
    A, B : in std_logic_vector(3 downto 0);
    Sum : out std_logic_vector(3 downto 0);
    done : out std_logic
  );
end serial_adder;

architecture behav of serial_adder is
  signal regA, regB, regSum : std_logic_vector(3 downto 0);
  signal carry : std_logic := '0';
  signal bit_cnt : integer range 0 to 4 := 0;
  signal working : boolean := false;
  signal s, c : std_logic;
begin
  s <= regA(0) xor regB(0) xor carry;
  c <= (regA(0) and regB(0)) or (regA(0) and carry) or (regB(0) and carry);

  process(clk, rst)
  begin
    if rst='1' then
      regA <= (others=>'0'); regB <= (others=>'0');
      regSum <= (others=>'0'); carry <= '0';
      bit_cnt <= 0; working <= false; done <= '0';
    elsif rising_edge(clk) then
      done <= '0';
      if start='1' and not working then
        regA <= A; regB <= B; regSum <= (others=>'0');
        carry <= '0'; bit_cnt <= 0; working <= true;
      elsif working then
        regA <= '0' & regA(3 downto 1);
        regB <= '0' & regB(3 downto 1);
        regSum <= s & regSum(3 downto 1);
        carry <= c;
        if bit_cnt=3 then
          working <= false; done <= '1';
        else bit_cnt <= bit_cnt+1; end if;
      end if;
    end if;
  end process;
  Sum <= regSum;
end behav;
\end{lstlisting}

\section{移位寄存器实现串行累加器}

\subsection{原理}

累加器 = 加法器 + 累加寄存器（ACC）。串行累加器用移位寄存器做ACC：每输入一个新操作数，将其与ACC中的值串行相加，结果存回ACC。

\textbf{结构：}1位全加器 + 进位DFF + ACC移位寄存器 + 输入移位寄存器。

\textbf{【VHDL核心逻辑——累加循环】}
\begin{lstlisting}
-- ACC累加：ACC = ACC + Input
-- 两个移位寄存器：acc_reg（累加值）和 in_reg（新输入）
-- 每拍：acc_reg(0) + in_reg(0) + carry → sum位存入acc高位，进位更新
...
if working then
  acc_reg <= sum_bit & acc_reg(7 downto 1);  -- 结果右移存入
  in_reg  <= '0' & in_reg(7 downto 1);        -- 输入右移
  carry   <= cout;
  if bit_cnt = 7 then working <= false; end if;
end if;
\end{lstlisting}

\section{移位寄存器实现乘法器}

\subsection{原理——移位相加法}

二进制乘法可以分解为移位+相加：

$$A \times B = A \times (B_3 \cdot 2^3 + B_2 \cdot 2^2 + B_1 \cdot 2^1 + B_0 \cdot 2^0)$$

逐位检查乘数B的每一位：
\begin{itemize}
  \item $B_0=1$：累加A（不移位）
  \item $B_1=1$：累加A左移1位（即$A\times2$）
  \item $B_2=1$：累加A左移2位
  \item ...
\end{itemize}

\textbf{更高效的方法（移位-加乘法器）：}每拍检查B的最低位，若为1则ACC += A，然后A左移1位，B右移1位。N拍后ACC中为乘积。

\textbf{逐拍分析：$5 \times 3 = 0101 \times 0011 = 15$}

\begin{center}
\begin{tabular}{|c|c|c|c|c|}
\hline
拍 & B(最低位) & 操作 & ACC & A(左移) \\
\hline
0 & 1 & ACC+=A(0101) & 0101 & 0101 \\
1 & 1 & ACC+=A(1010) & 1111 & 1010 \\
2 & 0 & 不加 & 1111 & 0100(移出) \\
3 & 0 & 不加 & 1111 & — \\
\hline
\end{tabular}
\end{center}

结果ACC=1111=15，正确。

\textbf{【VHDL——4位移位-加乘法器】}
\begin{lstlisting}
library ieee;
use ieee.std_logic_1164.all;
use ieee.numeric_std.all;

entity shift_multiplier is
  port(
    clk, rst, start : in std_logic;
    A_in, B_in : in unsigned(3 downto 0);
    product : out unsigned(7 downto 0);
    done : out std_logic
  );
end shift_multiplier;

architecture behav of shift_multiplier is
  signal A_reg : unsigned(7 downto 0);
  signal B_reg : unsigned(3 downto 0);
  signal ACC : unsigned(7 downto 0) := (others=>'0');
  signal cnt : integer range 0 to 4 := 0;
  signal working : boolean := false;
begin
  process(clk, rst)
  begin
    if rst='1' then
      ACC<=(others=>'0'); cnt<=0; working<=false; done<='0';
    elsif rising_edge(clk) then
      done <= '0';
      if start='1' and not working then
        A_reg <= "0000" & A_in;
        B_reg <= B_in;
        ACC <= (others=>'0');
        cnt <= 0; working <= true;
      elsif working then
        if cnt=4 then
          working <= false; done <= '1';
        else
          if B_reg(0)='1' then ACC <= ACC + A_reg; end if;
          A_reg <= A_reg(6 downto 0) & '0';  -- 左移1位
          B_reg <= '0' & B_reg(3 downto 1);   -- 右移1位
          cnt <= cnt+1;
        end if;
      end if;
    end if;
  end process;
  product <= ACC;
end behav;
\end{lstlisting}

\section{考试常见变式}

\begin{enumerate}
  \item \textbf{变式1：实现串并转换（SIPO）}

  串行数据输入，4个时钟周期后，并行输出完整4位数据。

  \item \textbf{变式2：实现并串转换（PISO）}

  先load并行数据，然后4个时钟周期移位输出。

  \item \textbf{变式3：多位移位（桶形移位器）}

  一次移多位（不是每次只移1位）。需要更复杂的组合逻辑。

  \item \textbf{变式4：移位寄存器+反馈 = LFSR/环形计数器}

  将串行输出反馈回串行输入。这将在第七部分详细展开。

  \item \textbf{变式5：移位寄存器的Testbench设计}

  如何验证移位寄存器的正确性。
\end{enumerate}

\section{Testbench}

\begin{lstlisting}[caption={移位寄存器Testbench}]
library ieee;
use ieee.std_logic_1164.all;

entity shift_reg_tb is
end shift_reg_tb;

architecture test of shift_reg_tb is
  signal clk, rst_n, si, so : std_logic;
  signal po : std_logic_vector(3 downto 0);
  constant clk_period : time := 20 ns;

  component shift_reg_right is
    port (clk, rst_n, si : in std_logic;
          so : out std_logic; po : out std_logic_vector(3 downto 0));
  end component;
begin
  uut: shift_reg_right port map (clk, rst_n, si, so, po);

  clk_process: process
  begin
    clk <= '0'; wait for clk_period/2;
    clk <= '1'; wait for clk_period/2;
  end process;

  stimulus: process
  begin
    rst_n <= '0'; si <= '0'; wait for 30 ns;
    rst_n <= '1'; wait for 5 ns;

    -- 输入序列：1, 0, 1, 1
    si <= '1'; wait for clk_period;
    si <= '0'; wait for clk_period;
    si <= '1'; wait for clk_period;
    si <= '1'; wait for clk_period;

    -- 观察：po应该依次为 1000, 0100, 1010, 1101
    wait;
  end process;
end test;
\end{lstlisting}

\begin{examtip}[移位寄存器秒杀法]
\textbf{右移}：\texttt{reg <= si \& reg(N-1 downto 1)}

\textbf{左移}：\texttt{reg <= reg(N-2 downto 0) \& si}

\textbf{并行装载}：\texttt{if load='1' then reg <= data\_in;}

\textbf{串行输出}：最高位/最低位 = so

\textbf{并行输出}：reg的全部位 = po
\end{examtip}

\section{变式详解}
\chapter{移位寄存器——全部变式详解}

\section{变式1：串并转换（SIPO）}

\begin{detailbox}[完整教学]
\textbf{【功能】}串行输入→4拍后并行读取。UART/SPI接收器的基础。
\textbf{【VHDL】}\texttt{reg <= si \& reg(N-1 downto 1); po <= reg;}
\end{detailbox}


\section{变式2：并串转换（PISO）}

\begin{detailbox}[完整教学]
\textbf{【功能】}load并行数据→逐拍串行输出。UART/SPI发送器的基础。
\textbf{【VHDL】}
\begin{lstlisting}
library ieee;
use ieee.std_logic_1164.all;

entity piso_shift_reg is
  generic (N : integer := 8);
  port (
    clk     : in  std_logic;
    rst_n   : in  std_logic;
    load    : in  std_logic;
    data_in : in  std_logic_vector(N-1 downto 0);
    so      : out std_logic
  );
end piso_shift_reg;

architecture behavioral of piso_shift_reg is
  signal reg : std_logic_vector(N-1 downto 0);
begin
  process(clk, rst_n)
  begin
    if rst_n = '0' then
      reg <= (others => '0');
    elsif rising_edge(clk) then
      if load = '1' then
        reg <= data_in;
      else
        reg <= '0' & reg(N-1 downto 1);
      end if;
    end if;
  end process;
  so <= reg(0);
end behavioral;
\end{lstlisting}
\end{detailbox}


\section{变式3：双向移位}

\begin{detailbox}[完整教学]
\textbf{【VHDL】}
\begin{lstlisting}
library ieee;
use ieee.std_logic_1164.all;

entity bidir_shift_reg is
  generic (N : integer := 8);
  port (
    clk   : in  std_logic;
    rst_n : in  std_logic;
    dir   : in  std_logic;
    si_r  : in  std_logic;
    si_l  : in  std_logic;
    po    : out std_logic_vector(N-1 downto 0)
  );
end bidir_shift_reg;

architecture behavioral of bidir_shift_reg is
  signal reg : std_logic_vector(N-1 downto 0);
begin
  process(clk, rst_n)
  begin
    if rst_n = '0' then
      reg <= (others => '0');
    elsif rising_edge(clk) then
      if dir = '1' then
        reg <= si_r & reg(N-1 downto 1);
      else
        reg <= reg(N-2 downto 0) & si_l;
      end if;
    end if;
  end process;
  po <= reg;
end behavioral;
\end{lstlisting}
\end{detailbox}


\section{变式4：循环移位}

\begin{detailbox}[完整教学]
\textbf{【功能】}数据不丢失——最低位移回最高位。
\textbf{【VHDL】}\texttt{reg <= reg(0) \& reg(N-1 downto 1);}（右循环）
\textbf{【本质】}这就是环形计数器的基础！
\end{detailbox}


\section{变式5：算术移位（保持符号位）}

\begin{detailbox}[完整教学]
\textbf{【逻辑右移】}高位补0 → \texttt{'0' \& reg(N-1 downto 1)}
\textbf{【算术右移】}高位补符号位 → \texttt{reg(N-1) \& reg(N-1 downto 1)}
\textbf{【应用】}有符号数÷2操作。
\end{detailbox}
